\documentclass[11pt]{article}
\usepackage[margin=1in]{geometry}
\usepackage{amsmath,amssymb}
\usepackage[version=4]{mhchem}
\usepackage[T1]{fontenc}
\usepackage{lmodern}
\usepackage[authoryear,round]{natbib}
\usepackage{graphicx,booktabs}
\usepackage{url}
\usepackage[left,mathlines]{lineno}
\usepackage[hidelinks]{hyperref}
\setlength{\emergencystretch}{2em}
\setlength{\tabcolsep}{4pt}
\renewcommand{\thefigure}{S\arabic{figure}}
\renewcommand{\thetable}{S\arabic{table}}
\title{Supporting Information for\\Precession-phase preference of Dansgaard--Oeschger event occurrence after accounting for background climate}
\author{}
\date{}
\begin{document}
\maketitle
\linenumbers
\noindent This Supporting Information contains Texts S1--S7, Table S1, and Figures S1--S10. Table S1 reports the assigned speleothem event times and their sensitivity to event-definition settings; complete event inventories and numerical diagnostics remain in the accompanying research data. References to numbered main figures refer to the accompanying article. All absolute-age axes place older ages on the left and younger ages on the right.

\section*{Text S1. Event selection, timing, and observation intervals}

The speleothem composite combines MF events 6.1--6.16 with five older Sofular events, labeled 6.9--6.13 in the source and renumbered 6.17--6.21 here (Figure~\ref{fig:S02}; Table~\ref{tab:S1}). We use the overlap between the records to avoid counting the same event twice. MF 6.16 corresponds to an unnumbered Sofular feature near 175.2 kyr BP; Sofular 6.9 is a distinct, older transition (Figure~\ref{fig:S01}). Sofular supplies the older part of the composite, including the interval containing the 7.08-kyr MF data gap, so this gap does not interrupt the composite observation interval.

We assign an onset time to each selected DO-like transition using the archived proxy series \citep{fohlmeister2023mfData,held2024sofularData}. We first divide each series at gaps exceeding 0.5 kyr, then interpolate each continuous segment onto a 1-yr grid and smooth it with a Gaussian kernel of 50-yr standard deviation. Within $\pm400$ yr of each published event's approximate age, we locate the steepest change in the expected direction. Measured toward older BP ages, this is a decrease in MF $\delta^{18}$O or an increase in Sofular $\delta^{13}$C. Search windows end at data gaps or at the midpoint between neighboring event labels, whichever is reached first. Table~\ref{tab:S1} lists the assigned times, and Figure~\ref{fig:S02} shows them on the proxy records. To test sensitivity to this procedure, we combine three smoothing scales (35, 50, and 75 yr) with three search half-windows (300, 400, and 500 yr). These nine settings change onset times but retain the same event identities.

The nominal observation windows are 12--123 kyr BP for the undisturbed NGRIP record \citep{ngrip2004high} and 132.5--204.5 kyr BP for the speleothem composite \citep{held2024dansgaard}. In each sequence, the oldest event initializes history, and the fitted interval runs from that event toward the younger observation boundary. The resulting intervals are 12--115.320 and 132.5--194.238 kyr BP, totaling 165.058 kyr. They include the time between the youngest warming and the younger boundary, but exclude the gap between the two records. These boundaries follow source coverage. We assume that the retained records adequately represent events within them; complete detection of every event has not been established. Descriptive event-phase plots retain all 55 events (Figure~2), including the two that initialize history but are excluded from the fitted event count.

Barker event identities and ages come directly from the warming-event columns of Supplementary Table~S3 in \citet{barker2011800}. We retain the published SpeleoAge ages and variable- and fixed-threshold classifications. Using the oldest event to initialize history leaves 69 fitted events over 0--396.464 kyr BP for the variable threshold and 58 over 0--392.246 kyr BP for the fixed threshold. Because both the fitted events and the interval lengths differ, their likelihoods cannot be used to rank the two definitions directly. The comparison tests event selection within one reconstruction; it does not provide two independent climate records.

\section*{Text S2. Event-age uncertainty and its effects on phase estimates}

We distinguish two sources of event-age uncertainty: where the onset is placed within a recorded transition (event-definition uncertainty), and the age assigned to that position by the chronology. The procedures below use the available error information to generate 10,000 alternative age sequences for each of NGRIP, the MF--Sofular composite, and Barker's variable-threshold catalogue. The same events remain in the same order in every sequence. Missing events are a separate issue, tested in Text~S7. All reported event ages use kyr before 1950 CE (kyr BP), after conversion from the source reference where needed.

\subsection*{NGRIP}

We perturb the 34 warming and 35 cooling boundaries together and preserve their alternating order, then retain the warming ages for the event-rate analysis. Definition uncertainties follow Table~2 of \citet{rasmussen2014stratigraphic} and range from a few years to about 200 yr. For each boundary, we draw an independent Gaussian age offset with mean zero and standard deviation given by its definition uncertainty (Figure~\ref{fig:S03}a).

Chronology errors are treated differently in the counted and modeled sections. Within the 12--123-kyr BP study window, annual layer counting extends to 60.152 kyr BP; older ages come from the model extension. The GICC05 maximum counting error (MCE) reaches 2.611 kyr at the counting endpoint. Following its interpretation as approximately $2\sigma$, we use MCE/2 as the chronology standard deviation \citep{svensson2008sixty,rasmussen2022gicc05}. The published GICC05modelext chronology has no quantified age uncertainty \citep{corrick2020synchronous}. \citet{moseley2020nalps19} extrapolated the counted-age uncertainty trend and obtained an uncertainty of about 4.5\% by 120 ka. Guided by that estimate, we use $\pm4.5\%$ of b2k age as an approximate $2\sigma$ envelope throughout the modeled section. The constant percentage across this section is our simplifying assumption.

A chronology error affects multiple events in the same record. To represent this dependence, we generate shifts on a 5-kyr grid from 0 to 120 ka b2k, spanning all retained boundaries. These are computational points, not additional dated horizons. We replace the 60-ka point with the exact counting endpoint at 60.202 ka b2k (60.152 kyr BP); the next point is 65 ka b2k. Each point receives the standard deviation described above, with zero uncertainty at age zero. Starting from the youngest point and moving toward older ages, we add independent Gaussian increments. If $V_j$ is the squared standard deviation at point $j$ and $V_{-1}=0$, these increments and the accumulated shifts are
\begin{equation}
 e_j\sim N(0,V_j-V_{j-1}),\qquad \delta_j=\sum_{m\leq j}e_m.
\end{equation}
Each older point therefore shares the younger points' accumulated error and adds a new increment. Before applying order constraints, its total variance is $V_j$. The standard deviation of each new increment is $\sqrt{V_j-V_{j-1}}$, so it is determined by the increase in variance, not by the difference between standard deviations.

We linearly interpolate the grid-point shifts to each event age, including across the counted--modeled boundary, and add both this chronology shift and the independent definition offset to the nominal age. If either the shifted grid ages or the combined warming/cooling boundaries change order, we reject the whole sequence and draw again. We retain 10,000 ordered sequences (Figure~\ref{fig:S03}b). The uncertainty envelopes determine the Gaussian standard deviations; they do not impose hard limits on the sampled shifts. Text~S5 tests alternative grid spacings.

\subsection*{MF and Sofular}

For the speleothems, event-definition uncertainty is represented by the nine smoothing and search-window settings in Text~S1. Table~\ref{tab:S1} lists the minimum and maximum onset shifts from the default setting for all 16 MF and five Sofular events. Most ranges are no wider than 40 yr, although MF events 6.10 and 6.14 span 211 and 365 yr. Each realization selects one of the nine settings with equal probability and applies it to all MF events; an independent selection is made for all Sofular events. Each record therefore uses one shared detection setting in a realization. We then add the chronology shifts described below, without adding further event-picking noise or changing event identities.

For MF, the archived stack provides nominal ages and lower and upper age bounds \citep{fohlmeister2023role,fohlmeister2023mfData}. We treat the bounds as an approximate 95\% uncertainty envelope and interpolate them within continuous proxy segments. At each assigned onset, we take the larger distance from the event age to either bound and divide it by 1.96 to obtain a Gaussian standard deviation. We then draw a single standard-normal value for each realization and multiply it by the local standard deviation at every MF event. The chronology component therefore shifts all MF events in the same direction, by amounts proportional to their local dating uncertainties. This shared factor accounts for the strong correlations in Figure~\ref{fig:S04}b; it is an assumed dependence structure, not a covariance estimated from the published age model.

For Sofular, the main ensemble uses ten So-4 U--Th controls and its published age--depth curve \citep{held2024dansgaard,held2024sofularData}. We locate each stack event at the depth on So-4 with the same reference age. At every dated depth, we draw an independent, zero-mean Gaussian offset with standard deviation equal to half the reported $2\sigma$ dating error. The offset is added to the published curve's age at that depth, not to the measured U--Th age. We linearly interpolate the two neighboring control offsets in the reference-age coordinate and add the result to the stack event age. This retains the curve's growth pattern between controls and recovers the published chronology when all offsets are zero. Shifted control ages must remain ordered. Interpolation does not cross the growth hiatus or extend beyond the dated or proxy coverage.

We combine the MF and Sofular draws and reject any proposal that reverses their composite event order. Figure~\ref{fig:S04}a shows the resulting event-age ranges. An alternative ensemble uses So-57 for source events 6.9--6.12 and retains So-4 for event 6.13, which lies outside So-57 coverage. We also shift the ages used to locate events on the individual stalagmite curves by up to $\pm1$ kyr. This tests the assumed location from which dating errors are transferred; the nominal stack event ages themselves are unchanged.

Locating stack events by age does not establish that their physical layers have been matched between stalagmites. Our procedure transfers dating errors through the published curves; it does not rerun StalAge or iscam \citep{scholz2011stalage,fohlmeister2012statistical}. It also excludes additional uncertainty from growth history, systematic dating errors, stack alignment, proxy measurements, and climatic lags between records.

\subsection*{Barker}

The Barker chronology ensemble uses the 70 variable-threshold events. Supplementary Table~S1 of \citet{barker2011800} provides 60 control points linking the synthetic record to the speleothem chronology (Figure~\ref{fig:S05}a). Their combined uncertainties include both matching error and absolute speleothem dating error. Because the source does not specify a probability distribution for these combined values, we use each as the half-width $U_j$ of a uniform distribution and independently draw $\delta_j\sim\mathrm{Uniform}(-U_j,U_j)$. We do not add the two component errors again. Event shifts are obtained by linear interpolation between the neighboring control shifts on the original SpeleoAge axis. For nominal event age $t$ between controls at $a_j$ and $a_{j+1}$,
\begin{equation}
 \delta(t)=(1-w)\delta_j+w\delta_{j+1},\qquad
 w=\frac{t-a_j}{a_{j+1}-a_j},
\end{equation}
and the realized event age is $t+\delta(t)$. Events between the same controls share their shifts, so their age errors are correlated.

Seven events lie in the alignment gap between controls at 264.24 and 317.70 kyr. We interpolate across this gap using those bounding controls. Four events are older than the oldest published control at 378.15 kyr, so we add a computational control at 400 kyr. Its uncertainty half-width is the larger of a linear extrapolation through the last four published uncertainties and the maximum of those four values, giving 1.30 kyr. We sample its shift in the same way as the other controls. Proposals are rejected if control ages change order; ordered controls also preserve the order of interpolated events. We retain 10,000 accepted sequences (Figure~\ref{fig:S05}b).

These combined control uncertainties describe chronological alignment and dating. They do not provide a separate picking error for each warming onset, so we add no further event-picking noise to the Barker chronology ensemble. Instead, we assess event-definition sensitivity by fitting the 59-event fixed-threshold catalogue at its nominal ages and comparing its phase response with the variable-threshold result (Main Figure~2; Text~S5). The definition comparison and chronology Monte Carlo analysis therefore address different sources of uncertainty.

\subsection*{Sensitivity of the chronology uncertainty}

For the primary catalogue, we randomly pair the 10,000 NGRIP realizations with the 10,000 speleothem realizations, using each once. A realization with any event outside its observation window is excluded without replacement or clipping. We retain 9,982 primary and all 10,000 Barker realizations. For each retained sequence, we refit BG and full, rebuilding history and starting each fitted interval at the realized oldest event. Climate and orbital records, nominal scaling, and younger observation boundaries remain fixed.

For NGRIP--MIS6, median $G$ is 0.2247 bits/event, with a 95\% age-sensitivity range of 0.1241--0.2914; nominal $p<0.05$ in 9,961 of 9,982 valid fits (99.79\%). For Barker varying threshold, median $G$ is 0.0986 bits/event, the corresponding range is 0.0708--0.1308, and 9,948 of 10,000 fits (99.48\%) meet the same threshold. The ranges are the 2.5th--97.5th percentiles across age realizations, not sampling confidence intervals; the threshold proportions are sensitivity summaries, not bootstrap $p$ values.

The chronology regions in Main Figure~3 use all 9,982 valid primary fits, and those in Figure~\ref{fig:barker_effect_precision} use all 10,000 Barker fits. Neither is based on the 200-realization subset used for the combined simulations. Text~S4 describes how these coefficient regions are constructed and converted to phase and rate-ratio bounds. Preferred phases remain near the precession-index minimum in both catalogues.

Using So-57 dating controls for the overlapping Sofular events (Text~S2) retains 9,982 of 10,000 age realizations within the observation windows. Median primary $G$ is 0.224 bits/event (central 95\% range 0.126--0.290), and preferred phase is $328.5^\circ$ ($311.3$--$342.7^\circ$). Nominal $p<0.05$ in 9,962/9,982 valid fits (99.80\%). These results closely resemble the main ensemble based on So-4.

We also vary the computational grid used to generate shared NGRIP chronology shifts. For grid spacings of 2.5, 5 and 10 kyr, separate 2,000-realization ensembles retain 1,993, 1,997 and 1,997 fits within the observation windows. Median $G$ is 0.223, 0.223 and 0.227 bits/event, and nominal $p<0.05$ in 99.75--99.95\% of retained fits. Each spacing is paired with the same speleothem realizations. These checks assess sensitivity to the assumed chronology correlations; their nominal tests are not calibrated by the main BG-model bootstrap.

\section*{Text S3. Model fitting and statistical tests}

We evaluate the likelihood integral by four-point Gauss--Legendre quadrature, splitting it at event ages and covariate interpolation knots. Event rates use history immediately before each event. Unlike the additive, self-exciting model of \citet{hawkes1971spectra}, our history term enters the log rate and is constrained to represent inhibition.

The phase multiplier is $\exp(\beta_s\sin\phi+\beta_c\cos\phi)$. Its maximum occurs at $\phi_{\max}=\operatorname{atan2}(\beta_s,\beta_c)\pmod{2\pi}$, and the maximum/minimum rate ratio is $\exp(2A)$, where $A=(\beta_s^2+\beta_c^2)^{1/2}$. These relationships also convert coefficient uncertainty into the effect ranges in Text~S4; preferred phase is undefined when $A=0$. The fitted counts in Main Figure~2a,c are obtained by integrating model rates over the time spent in each phase sector; the bars show observed phase distributions.

As a separate descriptive check, we apply the Rayleigh test for phase concentration against an independent, uniform-phase reference \citep{mardia2000directional}. Using all events, including those that initialize history, gives nominal $p=0.058$ for NGRIP--MIS6 ($n=55$), 0.079 for Barker varying threshold ($n=70$), and 0.047 for Barker fixed threshold ($n=59$). Only the last falls below 0.05. Limited event counts may reduce power, and pooling events across different climate states can obscure a conditional phase association. These are possible explanations for the weaker evidence, not causes established by our analysis. Unlike the model comparison, Rayleigh does not account for climate, history, or the time observed at each phase.

For bootstrap calibration, each catalogue uses its own fitted BG model and observation intervals. We simulate continuous event times by thinning, accepting candidate events according to the generating rate and updating history after each accepted event \citep{ogata1981simulation}. Event totals vary, and both models are refitted. With $B$ simulations and $b$ simulated $LR$ values at least as large as observed, $p_{\rm boot}=(b+1)/(B+1)$. All 9,999 simulations per catalogue yield valid fits. Text~S4 uses the same simulation procedure with the full model to assess effect precision. Note that for Barker's data, chronology and sensitivity analyses in Text S4 use the variable-threshold catalogue only.

\section*{Text S4. Event-sequence variability and combined uncertainty}

Even with known event ages, a finite sequence would not determine the preferred phase and rate modulation precisely: the same event-generating process can produce different sequences. We assess this sampling variability and its combination with age uncertainty for NGRIP--MIS6 (Main Figure~3) and Barker varying threshold (Figure~\ref{fig:barker_effect_precision}). Chronology uncertainty alone is described in Text~S2.

For sampling at fixed chronology, we generate 5,000 event sequences from each catalogue's nominal-age full model, retaining each segment's oldest conditioning event and younger boundary. Subsequent event times and totals vary. Each sequence is refitted with BG and full, re-estimating all coefficients.

For sampling combined with chronology uncertainty, we randomly select 200 distinct valid age realizations from each catalogue's ensemble and fit a full model to each. We simulate and refit 50 sequences per model, giving 10,000 estimates with equal weight per chronology. Each generator retains its realized oldest conditioning event and the original younger boundary; simulated events receive no additional age perturbation. Both analyses update history after every event and interpolate climate covariates at simulated ages, retaining nominal climate scaling (Text~S3). All 15,000 simulations per catalogue yield valid fits.

Each fit gives a sine and a cosine phase coefficient. To compare the three uncertainty sources consistently, we construct an ellipse for each set of coefficient pairs. Their spread and correlation determine its shape and orientation. We center each ellipse on the nominal-age estimate and enlarge it until it contains 95\% of the corresponding pairs. This retains shifts away from the nominal estimate, rather than recentering on the simulation mean. For sampling at fixed chronology, the ellipse is an approximate joint 95\% confidence region, conditional on the fitted model and chronology \citep{davison1997bootstrap}. The chronology and combined ellipses use the same construction but summarize sensitivity to the assumed age-error distributions; they are not calibrated confidence regions.

Panel c shows these regions. Considering all coefficient pairs within each ellipse gives the preferred-phase and rate-ratio bounds in panels a,b and the complete curve envelope in panel d. Thus, the curve limits are obtained from the regions as a whole, rather than from separate percentiles at each phase. If an ellipse includes $(0,0)$, it permits no phase effect and the preferred phase is unrestricted. Phase ranges continue across $360^\circ$ when needed.

For Barker, sampling gives a preferred-phase range of $282.2^\circ$--$389.1^\circ$ and a rate-ratio range of 1.25--8.84; combining sampling and chronology gives $281.9^\circ$--$390.3^\circ$ and 1.24--8.95 (Figure~\ref{fig:barker_effect_precision}). Both catalogues' sampling confidence regions exclude zero, but leave modulation strength imprecisely estimated. 

\section*{Text S5. Sensitivity to analysis choices}

We first test whether the phase association depends on how long earlier events influence the rate. Decay times of 1, 1.5, 2, 3 and 5 kyr are compared using the same starting events, 53 response events, and 165.058 kyr of fitted time. $G$ spans 0.243--0.252 bits/event and preferred phase $318.2$--$330.3^\circ$; all nominal LR $p$ values are below 0.00014. 

We then test the assumption that unobserved events before the starting event have no residual influence. This might be a problem for the MIS~6 speleothem, as it ends at approximately 194~kyr~BP, which is still during a glacial period, and there may have been events before it. At $\tau=1.5$ kyr, we add an initial history weight $h_{\rm pre}=0$, 0.5 or 1 to both segments. Its contribution decays as $h_{\rm pre}\exp(-u/\tau)$, in addition to the observed-event history, and all fitted time is retained. $G$ ranges from 0.248 to 0.250 and preferred phase from $327.8^\circ$ to $328.1^\circ$ (Figure~\ref{fig:S08}). These initial weights are alternative assumptions, not draws from an estimated uncertainty distribution.

We also assumed a linear relationship between the climate background and the log event rate. To examine the influence of a potential non-linear relationship, we add the square of LR04, the square of \ce{CO2}, both squares, or both squares plus the LR04--\ce{CO2} product, while retaining the linear terms. Each alternative is used in both BG and full, so their difference still tests only the added phase terms. $G$ spans 0.239--0.252 bits/event and preferred phase $318.0$--$328.1^\circ$. Adding LR04 squared lowers full-model AIC by 4.77 and shifts the phase maximum to $318.2^\circ$ (Figure~\ref{fig:S09}a--c).

Alternative history descriptions use the number of strictly earlier events in the preceding 1.5 kyr, time $w$ since the most recent earlier event, or $\log(1+w/1\ {\rm kyr})$. The exponential and event-count coefficients are constrained to be nonpositive; the two elapsed-time coefficients are unrestricted. All fits use the same events and fitted intervals, with history recalculated at event and integration times. For the exponential, count, elapsed-time, and log-elapsed-time versions, $G$ is 0.248, 0.219, 0.202 and 0.209 bits/event, respectively; preferred phases range from $328.1^\circ$ to $332.4^\circ$ (Figure~\ref{fig:S09}d--f). Separately, allowing NGRIP and the speleothem segment to have different phase coefficients gives nominal $p=0.810$ relative to the common-phase model. This provides no evidence for a difference, but does not establish that their responses are identical.

Barker's fixed-threshold catalogue tests sensitivity to the event-selection rule. At nominal ages it gives $G=0.107$ bits/event, nominal $p=0.0137$, bootstrap $p=0.0172$, preferred phase $347.1^\circ$ and rate ratio 3.19 (Main Figure~2). Its different oldest event also changes the fitted interval (Text~S1). Its separate BG-model bootstrap is described in Text~S3. 

We additionally compare transition directions within NGRIP, using the 35 stadial (GS) and 34 interstadial (GI) starts from \citet{rasmussen2014stratigraphic}. Each direction is fitted separately with the main continuous-time model, using LR04, CO$_2$ and same-direction exponential history ($\tau=1.5$ kyr, $\beta_H\leq0$) in BG and adding the two phase terms in full. Both use the 12--123-kyr BP observation window. Conditioning on GS-26 at 119.09 kyr BP or GI-25 at 115.32 kyr BP leaves 34 cooling or 33 warming response events and 107.09 or 103.32 kyr of fitted time, respectively. Each direction retains its own nominal climate scaling and receives 9,999 BG bootstrap simulations using the procedure in Text~S3. All bootstrap fits succeed. Cooling gives $G=0.136$ bits/event, phase maximum $340.1^\circ$, nominal $p=0.0408$ and bootstrap $p=0.0721$; warming gives $G=0.194$, $327.1^\circ$, nominal $p=0.0119$ and bootstrap $p=0.0298$. These are exploratory tests, without a multiple-testing adjustment, and the two directions are not independent replications.

For age sensitivity, both directions reuse the same 10,000 joint NGRIP chronology and definition-error realizations (Text~S2). Each valid fit rebuilds the oldest conditioning event, exposure, and history, retaining the nominal climate scaling. The cooling and warming analyses retain 9,263 and 9,982 draws, respectively; 737 and 18 are excluded because the oldest event crosses the 123~kyr~BP limit. None of the numerical fits fail. The nominal significance threshold $p<0.05$ is met in 5,433/9,263 cooling fits (58.7\%) and 8,035/9,982 warming fits (80.5\%). The median phase maxima are $332.5^\circ$ and $328.1^\circ$, with central 95\% age-sensitivity ranges of $308.0$--$350.9^\circ$ and $302.9$--$355.1^\circ$. The relatively low proportion of significant fits in NGRIP event sequence may be partly due to the limited number of draws, which reduces statistical power.

\section*{Text S6. Alternative orbital predictors}

These comparisons ask whether eccentricity, obliquity, or 65$^\circ$N summer-solstice insolation adds information about event occurrence, and whether the phase association remains after each is included. We fit eight models per catalogue: BG, BG plus phase (Pre), and BG plus each of the three alternative orbital variables, with and without Pre. BG retains climate and history terms and, for the primary catalogue, the segment term. Each alternative orbital variable is centered on its time-weighted mean and divided by its range over the nominal fitted intervals. Daily mean top-of-atmosphere insolation is calculated at solar longitude $90^\circ$ with a solar constant of 1,365 W m$^{-2}$. Its 1-kyr series is interpolated to event and integration times \citep{laskar2004long}.

For each alternative variable, three nested comparisons test its addition to BG, its addition when Pre is already included, and the addition of Pre when that variable is already included. The alternative variable adds one coefficient; Pre adds two (coefficient for sin and cos term). We apply Holm's correction to the nine nominal LR $p$ values within each catalogue to account for these multiple comparisons \citep{holm1979simple}. Figure~\ref{fig:orbital_comparison} and the main Results report the comparisons. 

For the orbital comparisons, we repeat the fits on 500 selected age realizations per catalogue. Of these, 499 primary and all 500 Barker realizations fall within the observation windows. Models compared within a realization share the same events, climate and orbital series, and fitted intervals. The resulting ranges describe age sensitivity; they do not provide additional bootstrap calibration under a null model.

\section*{Text S7. History contribution, model fit and additional event loss}

To test whether recent events help explain the rate after climate and precession phase are included, we compare the full model with a version that omits $H$. The null is $\beta_H=0$ and the allowed alternative is $\beta_H\leq0$, with $\tau=1.5$ kyr fixed. 

Zero lies at the edge of the allowed parameter range, so the usual $\chi^2_1$ reference need not apply \citep{self1987asymptotic}. 

We instead simulate 4,999 catalogues from each fitted no-history model, refit the models with and without history, and compare their LR distribution with the observed LR. Using the plus-one formula in Text~S3 gives bootstrap $p=0.026$ for the NGRIP-MIS6 ($LR=8.107$) and $p=0.012$ for Barker ($LR=7.722$). These tests support an inhibitory history contribution within the specified model family; they do not identify a particular physical recovery mechanism.







We next ask whether the full model adequately describes event spacing. An interval should tend to be short when the fitted rate is high and long when it is low. Time rescaling accounts for this variation by integrating the fitted rate between successive events \citep{brown2002rescaling}:
\begin{equation}
 z_{ri}=\int_{u_{r,i-1}}^{u_{ri}}\widehat\lambda_r(u)\,du,\qquad U_{ri}=1-\exp(-z_{ri}),\qquad u_{r0}=0.
\end{equation}
The quantity $z_{ri}$ measures interval length in units of accumulated fitted rate. For an uncensored sequence generated by a correctly specified conditional intensity with known parameters, time rescaling produces independent unit-exponential intervals; transforming them to $U_{ri}$ gives a uniform 0--1 scale. We use two prespecified checks: the Kolmogorov--Smirnov distance $D$ measures departure of the pooled $U$ distribution from uniformity, while $T_{\rm dep}$ measures association between neighboring transformed intervals:
\begin{equation}
 D=\sup_{0\leq x\leq1}|\widehat F_U(x)-x|,\qquad
 T_{\rm dep}=\frac{\left|\sum_r\sum_{i=1}^{N_r-1}(U_{ri}-0.5)(U_{r,i+1}-0.5)\right|}{\sqrt{\max(1,\sum_r\max(N_r-1,0))}}.
\end{equation}
Here $\widehat F_U$ is the empirical cumulative distribution of the transformed intervals. Neighboring pairs are formed within each segment and never cross record gaps. The interval from the youngest event to the younger observation boundary is incomplete: its rate integral is recorded as a censored endpoint, but it is not treated as a completed waiting time. Simulations with no response events have both statistics set to zero and remain in the calibration sample.

Both estimating the coefficients from the same events and ending observations at fixed boundaries affect the interval sample. We therefore calibrate these checks by simulation rather than applying the standard uniform-sample test directly. We simulate from the nominal-age full model, refit it, and recompute both statistics using the original observation boundaries. This follows the simulate--refit approach discussed by \citet{kling2026goodness}. For our inhibitory process it is an approximate parametric bootstrap; we do not invoke their theoretical guarantee for Poisson processes. The primary check reuses the 5,000 nominal-age simulations of Text~S4; Barker uses 1,999 full-model simulations. All replicates contribute, and no chronology perturbations are included. We apply Holm correction to the two diagnostic $p$ values within each catalogue. Adjusted values for $D$ and $T_{\rm dep}$ are 0.167 and 0.162 for the primary catalogue and 0.623 for both Barker statistics. Neither check detects clear misfit, although these two tests cannot establish that every feature of the event sequence is reproduced.

Finally, we test how further omissions from the event inventory would affect the phase estimate. We independently delete each event after the starting event with probability 0.1 or 0.2, retaining the original starting events and observation boundaries. For each deletion probability, we generate 500 random subsets, applying deletion to both primary segments, to MIS6 alone, or to Barker's variable-threshold catalogue. We refit BG and full after rebuilding history from the retained events; climate scaling stays fixed. All refits succeed. No age perturbations or further bootstrap simulations are added to these subsets. At 20\% deletion from both primary segments, the central 95\% range of phase shifts from the original estimate is $-14.7^\circ$ to $17.3^\circ$; for Barker it is $-19.2^\circ$ to $25.3^\circ$. These ranges describe sensitivity to additional random loss, not confidence intervals or estimates of which events are missing. They cannot exclude detection bias that depends on climate or phase.

\clearpage
\setcounter{table}{0}\refstepcounter{table}\label{tab:S1}
\section*{Table S1. Speleothem event ages and definition uncertainty}
\noindent Event identities follow \citet{fohlmeister2023role} for Melchsee--Frutt (MF) and \citet{held2024dansgaard} for Sofular. Ages are assigned using 50-yr Gaussian smoothing and a $\pm400$-yr search window (Text~S1), in kyr before 1950 CE; three decimal places specify numerical timing, not dating precision. Minimum and maximum offsets describe variation across nine smoothing/search-window combinations (Texts~S1--S2), with negative values indicating younger ages and positive values older ages. These ranges exclude chronology uncertainty and are not statistical confidence limits.

% Table S1 is maintained manually.
{\small
\centering
\begin{tabular}{lllrrr}
\toprule
Composite event & Record & Published label & Age (kyr BP) & Min. offset (yr) & Max. offset (yr) \\
\midrule
MIS 6.1 & MF & MF 6.1 & 134.976 & -6 & 7 \\
MIS 6.2 & MF & MF 6.2 & 138.393 & 0 & 8 \\
MIS 6.3 & MF & MF 6.3 & 146.251 & -14 & 21 \\
MIS 6.4 & MF & MF 6.4 & 146.847 & -11 & 8 \\
MIS 6.5 & MF & MF 6.5 & 148.439 & -4 & 5 \\
MIS 6.6 & MF & MF 6.6 & 149.458 & -13 & 13 \\
MIS 6.7 & MF & MF 6.7 & 150.462 & -15 & 0 \\
MIS 6.8 & MF & MF 6.8 & 151.115 & -4 & 8 \\
MIS 6.9 & MF & MF 6.9 & 159.522 & -7 & 5 \\
MIS 6.10 & MF & MF 6.10 & 160.879 & -205 & 6 \\
MIS 6.11 & MF & MF 6.11 & 165.006 & -28 & 1 \\
MIS 6.12 & MF & MF 6.12 & 168.566 & -15 & 4 \\
MIS 6.13 & MF & MF 6.13 & 169.675 & -1 & 5 \\
MIS 6.14 & MF & MF 6.14 & 171.506 & -365 & 0 \\
MIS 6.15 & MF & MF 6.15 & 172.364 & -5 & 2 \\
MIS 6.16 & MF & MF 6.16 & 174.752 & -1 & 1 \\
MIS 6.17 & Sofular & D-O 6.9 & 176.779 & -1 & 1 \\
MIS 6.18 & Sofular & D-O 6.10 & 179.831 & -23 & 9 \\
MIS 6.19 & Sofular & D-O 6.11 & 180.309 & -13 & 20 \\
MIS 6.20 & Sofular & D-O 6.12 & 190.972 & -21 & 19 \\
MIS 6.21 & Sofular & D-O 6.13 & 194.238 & -10 & 11 \\
\bottomrule
\end{tabular}
\par}
\clearpage
\begin{figure}[p]
\centering
\includegraphics[width=\linewidth,height=0.65\textheight,keepaspectratio]{figures/FigS01.pdf}
\caption{Sofular--Melchsee--Frutt (MF) correspondence across the overlapping interval. (a) Sofular $\delta^{13}$C and published labels \citep{held2024dansgaard}. (b) MF $\delta^{18}$O and published labels \citep{fohlmeister2023role}. Fine traces show observations; thick curves show Gaussian smoothing with $\sigma=50$ yr. Gray connectors mark visually inferred correspondences without tuning either chronology. MF 6.16 corresponds to an unnumbered Sofular feature near 175.2 kyr BP; Sofular 6.9 is a separate older event. Dashed leaders identify approximate label anchors, and gray shading marks the MF data gap. The Sofular proxy axis is reversed.}
\label{fig:S01}
\end{figure}
\clearpage
\begin{figure}[p]
\centering
\includegraphics[width=\linewidth,height=0.65\textheight,keepaspectratio]{figures/FigS02.pdf}
\caption{The 21-event MF--Sofular sequence. (a) Composite event times and source records, with the displayed division at 175 kyr BP. (b) MF $\delta^{18}$O and events 6.1--6.16 \citep{fohlmeister2023role}. (c) Sofular $\delta^{13}$C, with published labels 6.9--6.13 in brackets and composite labels 6.17--6.21 \citep{held2024dansgaard}. Gray points show observations; colored curves show smoothing with $\sigma=50$ yr. Filled circles mark maximum-gradient event times assigned near literature anchors (Text~S1). Open squares and horizontal bars show individual U--Th control ages and reported $2\sigma$ errors, rather than event-age intervals. The Sofular proxy axis is reversed.}
\label{fig:S02}
\end{figure}
\clearpage
\begin{figure}[p]
\centering
\includegraphics[width=\linewidth,height=0.65\textheight,keepaspectratio]{figures/FigS03.pdf}
\caption{NGRIP event-age uncertainty. (a) Working definition-error standard deviations for 34 warming boundaries (orange circles) and 35 cooling boundaries (blue squares) \citep{rasmussen2014stratigraphic}. The black line joins chronology envelopes: published maximum counting errors and the adopted 4.5\%-of-b2k-age model-extension envelope, both interpreted approximately as $2\sigma$ (Text~S2). (b) Combined chronology-plus-definition offsets: 30 realizations (gray), median (blue), and pointwise 2.5th--97.5th percentiles (shading) of 10,000 ordered realizations. Positive offsets indicate older ages. The dashed boundary separates counted and modeled ages at 60.152 kyr BP. These ranges describe the working uncertainty model.}
\label{fig:S03}
\end{figure}
\clearpage
\begin{figure}[p]
\centering
\includegraphics[width=\linewidth,height=0.65\textheight,keepaspectratio]{figures/FigS04.pdf}
\caption{Chronology and event-timing sensitivity of the speleothem composite. (a) Age offsets from nominal event times, with nominal ages beside labels. Dots mark medians; thick and thin bars span the 16th--84th and 2.5th--97.5th percentiles of 10,000 ordered realizations. Red and purple distinguish MF and Sofular; the dashed line marks zero offset. (b) Pearson correlations among sampled age anomalies, with the white boundary separating records. The distributions reflect the shared MF chronology factor, So-4 dated-depth perturbations, and alternative timing settings described in Text~S2; they are not full age-model posterior intervals.}
\label{fig:S04}
\end{figure}
\clearpage
\begin{figure}[p]
\centering
\includegraphics[width=\linewidth,height=0.65\textheight,keepaspectratio]{figures/FigS05.pdf}
\caption{Barker SpeleoAge chronology uncertainty \citep{barker2011800}. (a) Published combined control uncertainties (black circles), used as uniform proposal half-widths, and the auxiliary 400-kyr control (blue diamond; 1.30 kyr). The connecting line is an interpolated offset bound. (b) Twenty-five ordered offset maps (gray), median (blue), and pointwise 2.5th--97.5th percentiles (shading) of 10,000 maps. Positive offsets indicate older ages; red ticks locate the 70 events. Ochre shading marks the alignment gap, across which offsets are interpolated. The blue dotted line marks the last published control at 378.15 kyr; pale blue shading extends to the auxiliary control. Ranges describe the adopted chronology scenario (Text~S2).}
\label{fig:S05}
\end{figure}
\clearpage
\begin{figure}[p]
\centering
\includegraphics[width=\linewidth,height=0.65\textheight,keepaspectratio]{figures/FigS06.pdf}
\caption{Conditional fit contributions of precession phase and background climate. Solid bars show the gain from adding precession sine/cosine terms to a model containing LR04 and CO$_2$; hatched bars show the gain from adding LR04 and CO$_2$ to a model containing precession phase. Each comparison adds two coefficients and refits all coefficients on the same event sequence and observation intervals, retaining event history and, for NGRIP--MIS6, the segment term. Colors distinguish the three catalogues; Barker fixed threshold is an alternative event definition. Gains are measured in bits per fitted event at nominal ages, without uncertainty intervals; they are not additive shares of explained variability or measures of causal importance.}
\label{fig:conditional_contributions}
\end{figure}
\clearpage
\begin{figure}[p]
\centering
\includegraphics[width=\linewidth,height=0.68\textheight,keepaspectratio]{figures/FigS07.pdf}
\caption{Uncertainty in the Barker varying-threshold phase effect, using the same construction as Main Figure~3 (Text~S4). Chronology uses all 10,000 valid age fits; sampling uses 5,000 nominal-age full-model simulations; sampling + chronology uses 50 simulations from each of 200 age-specific full models. (a,b) Preferred-phase and rate-ratio bounds projected from the 95\% coefficient regions in (c); dots and dashed vertical lines mark nominal estimates. (c) Chronology (gray dotted), sampling (blue solid), and combined (orange dashed) regions, nominal estimate (black dot), and no phase effect (cross). (d) Envelopes from the complete regions, with the nominal multiplier in black and the sampling envelope shaded blue. Minimum/Maximum refer to the precession index. Only the sampling region is an approximate conditional confidence region; the other regions summarize sensitivity to the assumed age errors.}
\label{fig:barker_effect_precision}
\end{figure}
\clearpage
\begin{figure}[p]
\centering
\includegraphics[width=\linewidth,height=0.65\textheight,keepaspectratio]{figures/FigS08.pdf}
\caption{Sensitivity to history decay and initial history, using the same 53 fitted events and 165.058 kyr of observation time. (a--c) Decay times $\tau=1$, 1.5, 2, 3 or 5 kyr, with no contribution from unobserved events before each starting event. (d--f) Initial weights of 0, 0.5 or 1 for this unobserved history, with $\tau=1.5$ kyr. Columns show log-likelihood gain per event, $G$ (bits/event), preferred phase, and maximum/minimum phase rate ratio. Open markers and dashed lines identify the main settings. Starting events, younger boundaries, climate scaling, and the nonpositive history-coefficient constraint remain fixed; connected points show alternative settings, not uncertainty intervals (Text~S5).}
\label{fig:S08}
\end{figure}
\clearpage
\begin{figure}[p]
\centering
\includegraphics[width=\linewidth,height=0.65\textheight,keepaspectratio]{figures/FigS09.pdf}
\caption{Sensitivity to the form of the climate and history effects, using the same 53 fitted events and 165.058 kyr of observation time. Columns show log-likelihood gain per event, $G$ (bits/event), preferred phase, and maximum/minimum phase rate ratio. (a--c) Linear LR04/CO$_2$ terms with either square, both squares, or both squares plus the LR04--CO$_2$ product (``+ interaction''); exponential history is retained. (d--f) Exponential history, the event count in the preceding 1.5 kyr, time $w$ since the last event, and $\log(1+w/1\ {\rm kyr})$; linear climate terms are retained. Points show nominal-age estimates without uncertainty intervals, and dashed lines mark the main specification. For each setting, BG and full differ only in their phase terms (Text~S5).}
\label{fig:S09}
\end{figure}
\clearpage
\begin{figure}[p]
\centering
\includegraphics[width=\linewidth,height=0.65\textheight,keepaspectratio]{figures/FigS10.pdf}
\caption{Fit improvement from alternative orbital predictors for the primary catalogue (a--c) and Barker varying-threshold catalogue (d--f). BG contains climate and history terms, including the baseline difference between NGRIP and MIS~6; Pre adds the phase sine/cosine terms, and Orb is the orbital variable named in each row. Columns test Orb added to BG, Orb added after phase, and phase added after Orb. Circles show $G$ (bits/event) at nominal ages, ticks show medians across age realizations, and bars span their 2.5th--97.5th percentiles (499 valid primary fits; 500 Barker fits). Dashed lines mark phase added to BG alone. Orb adds one coefficient and Pre two; bars show age sensitivity rather than sampling confidence intervals. The nine additional tests per catalogue use Holm-adjusted nominal $p$ values (Text~S6).}
\label{fig:orbital_comparison}
\end{figure}
\clearpage
\bibliographystyle{agu}
\bibliography{reference}
\end{document}
